% Cover letter, compiled with XeLaTeX / Tectonic alongside resume/base/fonts (black-and-white)
% Upwork proposal: Technical Project Coordinator - Software Engineering & Delivery Management (invited; paste cover-letter.txt into the proposal box)
\documentclass[11pt]{article}
\usepackage[letterpaper,margin=1.1in,top=0.9in]{geometry}
\usepackage{fontspec,xcolor,enumitem}
\setmainfont[Path=fonts/lato/,BoldFont=Lato-Bol.ttf]{Lato-Reg.ttf}
\newfontfamily\namefont[Path=fonts/raleway/]{Raleway-Regular.otf}
\setlength{\parindent}{0pt}
\setlength{\parskip}{0.7em}
\setlist{nosep,leftmargin=1.4em}
\pagestyle{empty}
\begin{document}
\raggedright

{\namefont\fontsize{20pt}{24pt}\selectfont Jack Sturtevant}\\[-0.3em]
jack@sturt.io

\vspace{1em}

Hi,

Thanks for the invite. I'm a hands-on engineer who has spent most of the last seven years as the lead engineer on small teams. That has meant turning what a founder or client wants into a plan, splitting it into tasks, and keeping the work on track. This role is that same job with the coordination as the main focus.

Background: I've been a professional software engineer since 2019. I started at AWS and Applied Intuition. Since then I've been a two-time founder/CTO and a contract lead engineer, usually as the first or most senior engineer.

Coordinating teams and clients:
\begin{itemize}
\item Lancer (founder/CTO): I led a 7-person team building a freelance marketplace with escrowed payments, from planning through launch. 55 clients funded 264 jobs on it, and we won \$30K at Solana's Grizzlython.
\item Ralli (founder/CTO): I led the engineering team on a real-time social sports betting app (Next.js, NestJS, Postgres, WebSockets) and scaled us from part-time to full-time.
\item Banger (interim CTO): I took the founders' idea to a working trading platform in eight weeks. I worked directly with them on scope. The product led directly to a \$250K raise.
\item Coordinator (lead engineer, contract, 2026): I led engineering on the client's AI platform: LLM fine-tuning and serving, voice and video pipelines, CI/CD, and an API gateway for third-party apps.
\item Applied Intuition: I built the internal tool the company used to manage developers, teams and customers. I also built a support system that cut customer response time from over a day to under an hour.
\end{itemize}

Technologies: TypeScript, React/Next.js and Node (NestJS, tRPC) on the web, Python for AI and data work, Go, Postgres, Docker, Kubernetes, AWS and GCP. I work hands-on with LLM integrations, RAG, and workflow automation. I also used Java earlier in my career, though not recently.

My approach:
\begin{enumerate}
\item Requirements: I turn each request into a short written spec covering the goal, what's in and out of scope, acceptance criteria, and open questions. The client confirms it before any work starts, which catches most misunderstandings early.
\item Tasks and priorities: I break the spec into tasks a developer can finish in a day or two. Then I order them by risk and value, so the riskiest and most important work happens first.
\item Technical review: I read the code and the designs, so I can spot implementation problems and blockers early instead of finding out at the deadline.
\item Timelines: each milestone has a target date and a weekly written update covering what shipped, what's next, and what's at risk. If a date slips, the client hears it early, along with options such as cutting scope or moving the date.
\end{enumerate}

Availability and rate: I'm available now, part-time, and can scale up to 30+ hours a week. My rate is [RATE]/hr.

A few questions so I can understand the setup:
\begin{enumerate}
\item How many developers would I be coordinating, and are they in-house or contractors?
\item Will I be talking to clients directly on your Upwork account, or through my own?
\item What project-management tools do you use today (Jira, Linear, Trello, GitHub)?
\end{enumerate}

Happy to get on a call.

\vspace{0.4em}
Jack Sturtevant

\end{document}
